\documentclass[11pt]{article}
\usepackage[margin=1in]{geometry}
\usepackage{graphicx}
\usepackage{booktabs}
\usepackage{longtable}
\usepackage{array}
\usepackage{natbib}
\usepackage{hyperref}
\usepackage{xcolor}

\title{ Dual-Target Docking for Multi-Target Hit Discovery: A Reproducible Empirical Study }
\author{ SWW2026 Workshop Team }
\date{ 2026-02-05 }





\begin{document}
\maketitle

\begin{abstract}

This study evaluates whether dual-target reinforcement-learning generation
increases the probability of obtaining simultaneous hits on two protein targets
relative to single-target baselines.

We implement a fully reproducible two-stage pipeline: Stage 1 performs molecular
generation and scoring, while Stage 2 performs aggregation, statistical testing,
visualization, and manuscript assembly.

The experimental protocol enforces fairness by using matched molecular budgets,
fixed seed sets, shared hit definitions, and uniform post-processing rules.
Final conclusions are derived from pair-wise comparisons between
\texttt{double} and both baselines (\texttt{single\_a}, \texttt{single\_b})
with multiple-testing control.


\end{abstract}

\noindent\textbf{Keywords:} dual-target docking, molecular generation, FREED++, Boltz-2, reproducibility\\
\textbf{Experiment Run ID:} \texttt{ adhoc }


\section{ Introduction }

Проект проверяет гипотезу: оптимизация генерации молекул одновременно под две мишени
(\texttt{double}) увеличивает вероятность dual-hit относительно однобелковых режимов
\texttt{single\_a} и \texttt{single\_b}.

Эксперимент строится как воспроизводимый пайплайн:
подготовка таргетов, генерация FREED++, скоринг Boltz-2, агрегирование, статистика,
визуализация и отчёт.

Ключевой принцип дизайна эксперимента --- честное сравнение режимов при одинаковом
бюджете молекул, фиксированных сидах и одинаковых правилах hit-оценки.



\section{ Related Work }

Базовый генератор --- FREED++ (fragment-based RL) из экосистемы FFREED.
В текущем репозитории upstream reward path адаптирован под Uni-Dock с подготовкой
лигандов через RDKit+Meeko для устойчивого воспроизведения.

Для финального скоринга используется Boltz-2 через CLI, что позволяет работать с
единой схемой структуры+аффинности. Этот шаг выполняется после генерации и не
подменяет training reward внутри FREED++.

Многоцелевая оптимизация в литературе обычно сравнивается между скаляризациями
(\emph{worst-case}, гармоническое/геометрическое среднее, взвешенные формы)
и Парето-подходами. В данной работе основной фокус --- строгий baseline-сравнительный
эксперимент между \texttt{single\_a}, \texttt{single\_b} и \texttt{double}.



\section{ Methods }

Генерация выполняется в трёх режимах:
\texttt{single\_a}, \texttt{single\_b}, \texttt{double}.
Для \texttt{double} используется worst-case reward на z-нормированных прокси-скорах.

Скоринг выполняется скриптом \texttt{scripts/score.py} по двум мишеням пары с записью
\texttt{scored\_joined.csv}. Далее метрики и статистика рассчитываются в
\texttt{scripts/compute\_metrics.py} с bootstrap-процедурами и поправкой на множественные сравнения.

Оркестрация разделена на два этапа:
\begin{enumerate}
\item Stage 1: \texttt{scripts/run\_full\_experiment.py} (generate + score).
\item Stage 2: \texttt{scripts/run\_post\_pipeline.py} (aggregate, metrics, plots, render, summary, LaTeX).
\end{enumerate}



\section{ Experimental Protocol and Fairness Controls }

All experiments are executed under a two-stage orchestration model.

\textbf{Stage 1} (\texttt{scripts/run\_full\_experiment.py}) performs generation and
scoring for each pair, mode, and seed.

\textbf{Stage 2} (\texttt{scripts/run\_post\_pipeline.py}) performs aggregation,
metric computation, statistical testing, figure generation, and manuscript assembly.

Fairness controls:
\begin{itemize}
\item identical seed set across modes;
\item identical molecule budget per run;
\item identical hit thresholds per pair from \texttt{configs/metrics.yaml};
\item identical score convention (lower is better);
\item run-level metadata and timestamped logs for traceability.
\end{itemize}

Any interrupted run is resumed only if required artifacts are missing; otherwise it is skipped.



\section{ Result Scope and Decision Criteria }

This manuscript distinguishes between:
\begin{itemize}
\item \textbf{operational results} (all generated artifacts and intermediate statistics),
\item \textbf{claim-bearing results} (tests that pass predefined significance and effect-size gates).
\end{itemize}

The primary claim criterion is improvement in dual-hit rate for \texttt{double}
against both single-target baselines on the majority of target pairs, with
FDR-adjusted significance support.

Secondary criteria include balance between target scores and practical robustness
across seeds.

Where data are incomplete, the manuscript records the gap explicitly instead of
backfilling with extrapolations.



\section{ Limitations }

Current limitations:
\begin{itemize}
\item computational cost remains high for full multi-pair, multi-seed runs;
\item final external scorer latency constrains rapid iteration;
\item conclusions depend on selected target pairs and may not generalize
to all polypharmacology regimes.
\end{itemize}

Methodological limitations:
\begin{itemize}
\item reward shaping and normalization choices can bias search trajectories;
\item docking-driven objectives remain proxy objectives rather than direct wet-lab endpoints.
\end{itemize}

These limitations are tracked to prevent over-claiming and to guide future protocol upgrades.



\section{ Conclusion and Next Steps }

The repository is organized to produce publication-ready artifacts from executable
pipeline stages, rather than from manual post-hoc editing.

Next steps for A*/Q1 readiness:
\begin{enumerate}
\item complete all planned pair/mode/seed runs without missing artifacts;
\item freeze software/data snapshots and rerun Stage 2 for final tables and figures;
\item tighten narrative around hypothesis acceptance/rejection using only pre-registered criteria;
\item prepare a compact supplementary section with reproducibility commands and checksums.
\end{enumerate}

This structure keeps scientific claims auditable and directly linked to generated evidence.




\section{Quantitative Results}

No metric table available. Run \texttt{scripts/compute\_metrics.py} first.



No statistical table available. Run \texttt{scripts/compute\_metrics.py} first.


\section{Figures}

No figures available. Run \texttt{scripts/plot\_results.py} and \texttt{scripts/render\_molecules.py}.


\section{Reproducibility and Artifact Traceability}
\subsection{Software Snapshot}
\begin{itemize}
\item Master config version: \texttt{ 2.4 }

\item Docker image digest: \texttt{ sha256:84621724917265ea8f72445ced88af441730b45342b4613f4355ed5663c71d4a }
\item Boltz version: \texttt{ 2.2.1 }
\item Docking program/version: \texttt{ unidock_gpu } / \texttt{ Uni-Dock v1.1.3 }
\end{itemize}

\subsection{Data Snapshot}

\begin{itemize}
\item Calibration set: \texttt{ data/calibration/zinc_druglike_calibration_1000_v2024.smi }
\item Baseline set: \texttt{ data/baseline/zinc_druglike_baseline_1000_seed42.smi }
\end{itemize}

\appendix
\section{Machine Summary Excerpt}

No summary excerpt available. Run \texttt{scripts/summarize.py}.


\bibliographystyle{plain}
\bibliography{references}

\end{document}